\documentclass[../../main.tex]{subfiles}
\begin{document}

\begin{flushright}
\footnotesize\textit{【自動文字起こし・要確認】}
\end{flushright}

{\bf [解]}
題意の数列
\begin{align}
a(n)=\dfrac{(n+2)(n+3)(n+4)}{n!} \label{eq:1}
\end{align}
について考える．

\paragraph{(1)}

$n\to\infty$の時を考えるので$n \ge 4$として良い．この時$n! \ge n(n-1)(n-2)(n-3)$だから\cref{eq:1}に代入して
\begin{align}
 0 \le a(n) \le \dfrac{(n+2)(n+3)(n+4)}{n(n-1)(n-2)(n-3)} = \dfrac{(1+2/n)(1+3/n)(1+4/n)}{n(1-1/n)(1-2/n)(1-3/n)}
\end{align}
である．右辺は$n\to\infty$で$0$に収束するから，挟み撃ちの定理より
\begin{align}
 \lim_{n\to\infty}a(n) = 0
\end{align}
である．

\paragraph{(2)}
$n=1$から$a(n)$を書き出すと
\begin{align}
\begin{split}
 a(1)&=a(2)=60 \\
a(3)&=35,\quad a(4)=14,\quad a(5)=\dfrac{21}{5},\quad a(6)=1 \\
a(7)&=\dfrac{11}{56}<1
\end{split} \label{eq:2}
\end{align}
である．

以下，$a(n)$が$n\ge2$で単調減少であることを示す．そのために$a(n)-a(n+1)$を計算すると
\begin{align}
\begin{split}
a(n)-a(n+1)
 &=\dfrac{(n+2)(n+3)(n+4)}{n!} - \dfrac{(n+3)(n+4)(n+5)}{(n+1)!} \\
 &=\dfrac{(n+3)(n+4)}{(n+1)!}\left[(n+1)(n+2)-(n+5)\right] \\
 &=\dfrac{(n+3)(n+4)}{(n+1)!}(n^2+2n-3)>0 \quad(\because n\ge2)
\end{split} 
\end{align}
より示された．これと$a(7)<1$，および(1)の結果$a_n\to0\ (n\to\infty)$より，$7\le n$なる$n\in\mathbb{N}$に対しては$0<a(n)<1$となり$a(n)$が整数になることはない．以上から$a(n)$が整数となる$n$は
\begin{align}
n=1,2,3,4,6
\end{align}
である．

\paragraph{(3)}
(2)で$n\ge7$の時に$0<a(n)<1$だったから$S_n=\displaystyle\prod_{k=1}^na(k)$とおくと，$n\ge7$の時には$S_n$は単調減少である．
また\cref{eq:2}から$n=1,2,\cdots,7$の時は$S_n\in\mathbb{Z}$である．そこで$n\ge8$の時を順番に計算すると
\begin{align}
\begin{split}
S(7)&=a(1)\cdots a(7)=2^2\cdot3^3\cdot5^2\cdot7^2\cdot11=1455300\in\mathbb{Z} \\
S(8)&=S_7\cdot a(8)=\dfrac{3^2\cdot5^2\cdot7\cdot11^2}{2^2}\notin\mathbb{Z} \\
S(9)&=S(8)\cdot a(9)=\dfrac{5\cdot11^3\cdot13}{2^7\cdot3}\notin\mathbb{Z} \\
S(10)&=S(9)\cdot a(10)=\dfrac{11^3\cdot13^2}{2^{12}\cdot3^4\cdot5}\notin\mathbb{Z}<1
\end{split}
\end{align}
である．従って$10\le n$の時は$0<S(n)<1$で$S(n)$は整数とならない．以上から求める整数は
\begin{align}
n=1,2,3,4,5,6,7
\end{align}
である．

\newpage

\end{document}
